\setAuthor{Kristian Kuppart}
\setRound{piirkonnavoor}
\setYear{2019}
\setNumber{G 3}
\setDifficulty{3}
\setTopic{Optika}

\prob{Kärbes}
Kärbes asub kumerläätsest kümnekordse fookuskauguse kaugusel ning läätse optilisest peateljest kolme fookuskauguse kaugusel ning hakkab liikuma otse oma kujutise poole. Kui kaugel läätse tasandist asub kärbes hetkel, kui tema tõeline kujutis liigub kärbse suhtes\\
\osa kõige aeglasemalt;\\
\osa kõige kiiremini? Kui suur on kärbse kujutise kiirus kärbse suhtes nendel hetkedel? Põhjendage vastust. 

\hint
Kärbse liikumise käigus on ta suunatud läätse keskpunkti suunas ning seega asuvad kärbes ja tema kujutis alati samal sirgel. Kärbse kujutis algab liikumist fookusest seespol ning mingil hetkel jõuab lõpmatusse.\solu
\begin{center}
  \includegraphics[width=0.9\textwidth]{2019-v2g-03-sol.pdf}
\end{center}

Olgu läätse fookuskaugus $f$, kärbse esialgne asukoht $A$, vähima kiirusega asukoht $B$, suurima kiirusega asukoht $C$ ning nende vastavate kujutiste asukohad tähistatud ülakomaga. Paneme tähele, et kuna kärbes lendab alati oma kujutise poole, on ta suunatud läätse keskpunkti suunas ning seega asuvad kärbes ja tema kujutis alati sirgel $AA'$.\\
\osa Kui kärbes lendab punktist $A$ punkti $B$, siis tema kujutis liigub punktist $A'$ punkti $B'$. Jooniselt on näha, et kärbse kujutis liigub aeglasemalt kui kärbes. Kui kärbes on läätsest kaugusel $2f$, siis asub samuti ta kujutis kaugusel $2f$ teisel pool läätse (selles saab veenduda läätsevalemi kaudu). Sellisel juhul on kärbse ja tema kujutise kiirused võrdsed ning suhteline kiirus seega null. Teisisõnu on kärbse ja tema kujutise vaheline kiirus minimaalne siis, kui kärbes asub läätsest kahekordse fookuskauguse kaugusel.\\
\osa Kui kärbes liigub punktist $B$ fookaaltasandi suunas, siis kujutise kiirus järjest suureneb ning vahetult enne fokaaltasandile jõudmist on kujutise kiirus lõpmatult suur ($v_\mathrm{max} \rightarrow \infty\SI{}{\,m/s}$). Seega on kujutise kiirus kärbse suhtes maksimaalne siis, kui kärbes on väga lähedal fokaaltasandile ehk kärbes asub läätse tasandist fookuskauguse kaugusel.\probend